% Auto-generated by scripts/run_monthly_hs300_comparison.py
% Source: monthly comparison generator output
% Do not edit by hand.
\begin{figure}[H]
\centering
\includegraphics[width=\textwidth]{improved_rrp_vs_hs300_monthly_comparison.png}
\caption{Improved Convex Adaptive Global RRP 与沪深300ETF的月度表现对比}
\label{fig:monthly_hs300_comparison}
\end{figure}

图~\ref{fig:monthly_hs300_comparison}展示了 2018-01 至 2026-07 期间 Improved Convex Adaptive Global RRP 与沪深300ETF的月度净值和月度超额收益。Improved Convex Adaptive Global RRP 累计收益为 59.97\%，沪深300ETF为 66.50\%；但策略月度波动率仅 0.79\%，低于沪深300ETF的 4.63\%。日频最大回撤分别为 -3.81\%（2020-03-19）和 -38.85\%（2024-02-02）。
策略在 53/103 个月跑赢沪深300ETF，月度相关系数为 0.565。该结果说明，本文模型并非以捕捉A股权益贝塔为目标；其优势主要体现为显著压低波动和回撤，在权益强上涨月份往往牺牲一定进攻弹性。
